\problem{6.1}{Parity, mirror reflections, and hydrogenic harmonics}
\textbf{(a)} A reflection in the $xy$ plane sends
$(x,y,z)\mapsto(x,y,-z)$. Following it with a
half turn about $z$ sends
$(x,y,-z)\mapsto(-x,-y,-z)=-\mathbf r$.
Thus spatial inversion can be realized by a
mirror reflection followed by a rotation.
\textbf{(b)} In spherical coordinates
$x=r\sin\theta\cos\phi$, $y=r\sin\theta\sin\phi$,
$z=r\cos\theta$. The transformation
\[
\theta\longmapsto\pi-\theta,\qquad
\phi\longmapsto\phi+\pi
\]
changes each Cartesian coordinate's sign,
and $r$ stays fixed. Consequently
\begin{equation}
\boxed{(\PP\psi)(r,\theta,\phi)=
\psi(r,\pi-\theta,\phi+\pi).}
\end{equation}
\textbf{(c)} For $\psi_{n\ell m}=R_{n\ell}(r)
Y_\ell^m(\theta,\phi)$, use the standard
harmonic $Y_\ell^m=N_{\ell m}
P_\ell^m(\cos\theta)\ee^{\ii m\phi}$.
Since
$P_\ell^m(-u)=(-1)^{\ell+m}P_\ell^m(u)$ and
$\ee^{\ii m\pi}=(-1)^m$,
\[
Y_\ell^m(\pi-\theta,\phi+\pi)
=(-1)^{\ell+2m}Y_\ell^m(\theta,\phi)
=(-1)^\ell Y_\ell^m.
\]
The radial function depends only on $r$:
\begin{equation}
\boxed{\PP\psi_{n\ell m}=(-1)^\ell\psi_{n\ell m}.}
\end{equation}
This conclusion is valid for \emph{any} central
potential, not just the Coulomb potential.
\phys{Even $\ell$ gives even parity; odd $\ell$
gives odd parity.}

\problem{6.2}{Exponential of a Hermitian generator is unitary}
Let $Q^\dagger=Q$ and
$U=\exp(\ii Q)=\sum_{n\ge0}(\ii Q)^n/n!$.
Taking the adjoint termwise,
\[
U^\dagger=\sum_{n\ge0}
\frac{(-\ii Q^\dagger)^n}{n!}
=\exp(-\ii Q).
\]
The two exponentials commute because
they contain the same operator $Q$.
Therefore
\begin{equation}
\boxed{U^\dagger U=
\ee^{-\ii Q}\ee^{\ii Q}=
\ee^{0}=I,\qquad UU^\dagger=I.}
\end{equation}
By the spectral theorem this remains
valid for a self-adjoint (possibly
unbounded) $Q$ through the functional
calculus; a formal power series need
not converge on every state.

\problem{6.3}{Momentum does not change under spatial translations}
Take $\T(a)f(x)=f(x-a)$,
$\p=-\ii\hb\partial_x$, and
$\T(a)^\dagger=\T(-a)$. Acting on
a differentiable test function,
\begin{align}
(\T^\dagger\p\T f)(x)
&=\T(-a)[-\ii\hb f'(x-a)]\\
&=-\ii\hb f'(x)=(\p f)(x).
\end{align}
Hence
\begin{equation}
\boxed{\T(a)^\dagger\p\T(a)=\p.}
\end{equation}
Equivalently $[\p,\T(a)]=0$
because $\T(a)=\exp(-\ii a\p/\hb)$.
\phys{A change of origin shifts position
but not the momentum.}

\problem{6.4}{Translation of an arbitrary operator series}
First, the example discussed in the text gives
$\T^\dagger\x\T=\x+a$; Problem 6.3 gives
$\T^\dagger\p\T=\p$.
For any operators $A,B$ and unitary $T$,
\[
T^\dagger(AB)T=
(T^\dagger AT)(T^\dagger BT)
\]
by inserting $TT^\dagger=I$.
Apply the identity repeatedly:
\[
\T^\dagger \x^m\p^n\T
=(\x+a)^m\p^n.
\]
Therefore if $Q=\sum_{m,n}c_{mn}\x^m\p^n$
in the ordering specified in the exercise,
\begin{equation}
\boxed{\T^\dagger Q(\x,\p)\T
=Q(\x+a,\p).}
\end{equation}
The result extends to ordered products
and suitable functions by functional
calculus, but operator ordering should
be preserved when $[\x,\p]\ne0$.

\problem{6.5}{Proving the periodic factor in Bloch's theorem}
Let the Hamiltonian commute with translation
by $a$, and let its chosen eigenstate obey
\[
\T(a)\psi(x)=\psi(x-a)
=\ee^{-\ii qa}\psi(x).
\]
Define $u(x)=\ee^{-\ii qx}\psi(x)$,
so $\psi(x)=\ee^{\ii qx}u(x)$.
Then
\begin{align}
u(x-a)
&=\ee^{-\ii q(x-a)}\psi(x-a)\\
&=\ee^{-\ii qx}\ee^{\ii qa}
\ee^{-\ii qa}\psi(x)
=u(x).
\end{align}
Shifting $x\mapsto x+a$ shows
$u(x+a)=u(x)$:
\begin{equation}
\boxed{\psi_q(x)=\ee^{\ii qx}u_q(x),
\qquad u_q(x+a)=u_q(x).}
\end{equation}
Here $q$ and $q+2\pi/a$ describe
the same Bloch translation eigenvalue.
A periodic potential conserves \emph{crystal}
momentum, not necessarily ordinary momentum.

\problem{6.6}{The lowest two bands in the periodic finite well}
\textbf{(a) Derivation.}
Write $\psi_q(x)=\ee^{\ii qx}u_q(x)$, where
$u_q(x+a)=u_q(x)$. Differentiating twice,
\[
\psi_q''=\ee^{\ii qx}
(u_q''+2\ii q u_q'-q^2u_q).
\]
Substitute in $-\hb^2\psi_q''/(2m)+V\psi_q=E\psi_q$:
\begin{equation}
\boxed{\left[-\frac{\hb^2}{2m}\frac{\dd^2}{\dd x^2}
-\frac{\ii\hb^2q}{m}\frac{\dd}{\dd x}
+\frac{\hb^2q^2}{2m}+V(x)\right]u_q
=E_q u_q.}
\end{equation}
The opposite Bloch-phase convention changes
the sign of the linear-derivative term but
not the energies.
\textbf{(b) Exact matching reference.}
Use $E_*=\hb^2/(2ma^2)$,
$s=x/a$, $\epsilon=E/E_*$, and
$v=-20$ on $1/4<s<3/4$ (zero outside).
Translation of the origin gives two
segments of length $a/2$ with potentials
$-20E_*$ and $0$. At a constant-potential
segment of dimensionless length $d$,
a state $(\psi,\psi')^T$ propagates through
\begin{equation}
M(k,d)=
\begin{pmatrix}
\cos(kd)&\sin(kd)/k\\
-k\sin(kd)&\cos(kd)
\end{pmatrix},
\qquad k=\sqrt{\epsilon-v}.
\end{equation}
For imaginary $k$ the entries are real
hyperbolic functions. Bloch boundary
conditions imply that the eigenvalues of
$M(\sqrt\epsilon,\tfrac12)
M(\sqrt{\epsilon+20},\tfrac12)$
are $\ee^{\pm\ii qa}$. Since $\det M=1$,
\begin{equation}
\boxed{\cos(qa)=\frac12\operatorname{tr}
\left[M(\sqrt\epsilon,\tfrac12)
M(\sqrt{\epsilon+20},\tfrac12)\right].}
\end{equation}
A root search gives the \emph{two lowest}
energy levels in units of $E_*$:
\begin{equation}
\begin{array}{c|rr}
qa& E_1/E_*&E_2/E_*\\\hline
-\pi&-6.981932&5.618248\\
-\pi/2&-10.060437&13.552029\\
0&-11.989299&28.848077\\
\pi/2&-10.060437&13.552029\\
\pi&-6.981932&5.618248
\end{array}
\end{equation}
\textbf{(c) Reproducible finite differences.}
On an $N$-point periodic grid with
$h=1/N$, approximate the Bloch-covariant
kinetic operator by
\begin{equation}
(H_qu)_j=\frac{
2u_j-\ee^{\ii qh}u_{j+1}
-\ee^{-\ii qh}u_{j-1}}{h^2}
+v(s_j)u_j,
\quad j\bmod N.
\end{equation}
The resulting matrix is Hermitian.
For $N=400$ or $800$, diagonalize
its two lowest eigenvalues for
$qa\in[-\pi,\pi]$; increasing $N$
checks convergence to the exact-matching
table. The plotting script writes
\texttt{figures/bloch\_bands.pdf}.
\phys{The two bands are even in $q$ because
the potential and Hamiltonian are real;
their band gap is an effect of periodic
scattering, not an imposed boundary wall.}

\problem{6.7}{Total momentum conservation for an interacting pair}
Let
$H=\p_1^2/(2m_1)+\p_2^2/(2m_2)+
V(|x_1-x_2|)$.
\textbf{(a)} A simultaneous translation
of both coordinates acts as
\begin{equation}
\T(a)\psi(x_1,x_2)
=\psi(x_1-a,x_2-a)
=\exp[-a(\partial_{x_1}+\partial_{x_2})]\psi.
\end{equation}
Using $p_j=-\ii\hb\partial_{x_j}$,
\[
\boxed{\T(a)=\exp(-\ii aP/\hb),\qquad
P=p_1+p_2.}
\]
\textbf{(b)} The kinetic energy
commutes with $P$.
With $r=x_1-x_2$,
\[
[P,V(|r|)]=-\ii\hb
(\partial_{x_1}+\partial_{x_2})V(|r|)
=0
\]
since $\partial_{x_1}V=V'(r)$ and
$\partial_{x_2}V=-V'(r)$ wherever
the derivative is defined. Therefore
\begin{equation}
\boxed{[H,P]=0,\qquad
\frac{\dd\av P}{\dd t}=0.}
\end{equation}
\phys{Internal forces can exchange
momentum between the particles but
cannot change the total momentum.}
